\section{Talos Requirements}
\label{app:talos_requirements}

Talos requires two main inputs to execute the complete exploration
flow: a DNN workload and a pool of characterized hardware IPs. The
user can also provide design constraints and exploration parameters.

This appendix summarizes the information required by the current
implementation.

\subsection{Workload}

The workload is provided as an ONNX model. This model is used by
ZigZag during Level 1 to evaluate the different accelerator
architectures and find a suitable mapping.

Talos obtains the numerical representation of the workload from the
element data type (\texttt{elem\_type}) declared for each tensor in
the ONNX model. The activation, weight, and output bit
widths are used to configure the workload evaluated by ZigZag and are
preserved in the activity profile passed to Level~2.

Talos does not perform quantization. Therefore, the ONNX model must
already use the numerical format that is intended to be evaluated.

For every architecture passed to Level 2, Talos stores the activity
information obtained from ZigZag. For each layer, this profile
contains:

\begin{itemize}
    \item execution cycles;
    \item number of MAC operations;
    \item number of PEs used by the mapping;
    \item accesses to each memory level;
    \item operand precision and numerical format.
\end{itemize}

This activity profile is reused during Level 2. The workload is not
mapped again for every physical IP combination.


\subsection{IP Pool}

The hardware components available to Talos are provided using an IP
pool in YAML format. A pool contains the technology node and a list of
characterized IP blocks. Each entry describes one fixed configuration.
Parameterizable IPs are not stored as a single entry with free
parameters. Each configuration that has to be available to Level~2,
such as the INT8 and INT16 versions of a PE or the different depths of
a register file, is characterized and stored as a separate entry with
its own ID.

A simplified structure is shown below; the complete pools used in the
final evaluation, together with an annotated example of one entry, are
reproduced in Appendix~\ref{app:ip_pools}.

\begin{verbatim}
technology_nm: 65

ips:
  - ...
\end{verbatim}

The technology node is shared by the complete pool. It is also used by
the Level 1 calibration when CACTI is executed.

The main IP types used by the current implementation are:

\begin{itemize}
    \item processing elements;
    \item register files;
    \item global buffers.
\end{itemize}

A DRAM characterization can also be included. However, DRAM is treated
as part of the platform and is not explored as a Level 2 component.


\subsection{Common IP Information}

Every IP is identified by an ID and its component type. In addition,
Talos uses the following physical information:

\begin{itemize}
    \item area;
    \item delay;
    \item maximum frequency;
    \item throughput;
    \item power characterization.
\end{itemize}

Memory components can additionally provide:

\begin{itemize}
    \item capacity;
    \item bandwidth.
\end{itemize}

These fields are represented in the IP pool using values such as:

\begin{verbatim}
id: ...
type: ...
area: ...
throughput: ...
delay: ...
fmax_mhz: ...
capacity_bits: ...
bandwidth_bits: ...
\end{verbatim}

Not every field is required for every component type. For example,
capacity has no meaning for a PE, while it is required to determine
whether a memory can implement an abstract memory from Level 1.

Area, throughput, and delay are part of the generic IP
representation. Maximum frequency and power characterization become
necessary when the corresponding physical constraints or
workload-aware evaluations are used. In this case, the selected IPs
must provide compatible operating-point information.

Some of these fields may not be used by the default exploration
objectives, but they are preserved as additional architectural and
physical information. They can also be used when customizing the
evaluation module or implementing alternative exploration
heuristics.


\subsection{Power Characterization}

Level 2 requires an idle and active power characterization for the
selected IPs. The characterization contains:

\begin{itemize}
    \item characterization source;
    \item activity method;
    \item reference frequency;
    \item idle power;
    \item active power;
    \item voltage;
    \item temperature;
    \item process corner.
\end{itemize}

An example is:

\begin{verbatim}
power_model:
  source: ...
  activity_method: ...
  reference_frequency_mhz: 200.0
  p_idle_w: ...
  p_active_w: ...
  voltage_v: 1.2
  temperature_c: 25.0
  corner: ...
\end{verbatim}

The idle and active powers are given per IP instance.

Talos uses these values at their characterization point. It does not
scale the power values to another voltage or frequency.

All the IPs used in the same implementation must have a compatible
operating point. Their maximum frequency must also be high enough to
support the reference frequency used during the workload evaluation.

The maximum frequency is only used as a physical limit. Talos does not
automatically execute a workload faster when an IP has a higher
maximum frequency.


\subsection{Processing Elements}

A PE requires additional information about its compute capability.
The main value used during Level 2 is:

\begin{verbatim}
metadata:
  macs_per_cycle: ...
\end{verbatim}

This value is used to check whether the selected PE can sustain the
mapping produced by ZigZag.

A PE can also describe its numerical representation:

\begin{verbatim}
metadata:
  precision_bits: ...
  numeric_format: ...
\end{verbatim}

Talos compares the PE numerical format and precision with the input
activation and weight requirements obtained from the ONNX workload. A
mismatch is deliberately accepted because the user may want to
evaluate IPs with a different representation. However, Talos reports
a warning so that the mismatch is not introduced without the user's
knowledge. The candidate is therefore kept in the exploration.

Other information, such as the pipeline latency or dynamic energy per
MAC, can also be stored as metadata when available. These values are
not used by the default evaluator, but they remain available for
customizing the evaluation module or implementing additional
estimation and selection heuristics without changing the IP-pool
format.


\subsection{Memory Components}

Register files and global buffers require a capacity and bandwidth:

\begin{verbatim}
capacity_bits: ...
bandwidth_bits: ...
\end{verbatim}

These values are first used to determine whether the physical memory
can implement the memory requirements generated in Level 1.

For workload-aware evaluation, memory IPs also require:

\begin{verbatim}
metadata:
  accesses_per_cycle: ...
\end{verbatim}

This represents the maximum number of accesses that the memory can
provide in one cycle.

Talos compares this value with the memory traffic obtained from the
ZigZag mapping. If the selected memory cannot sustain the required
traffic, that physical implementation is rejected.

Since ZigZag provides aggregate memory accesses and the characterized
port rates are approximate, the current implementation accepts a
utilization tolerance of 20\%. Candidates above this tolerance are
rejected, while accepted utilization values are limited to one for
the power interpolation.


\subsection{Processing Elements with Included Register Files}

Some PE implementations already contain local memories as part of
their RTL.

Talos allows these register files to be declared as part of the PE:

\begin{verbatim}
included_rfs:
  rf_i1: <register-file IP>
  rf_i2: <register-file IP>
  rf_o:  <register-file IP>

included_rf_power_mode: parent_idle_baseline
\end{verbatim}

The referenced register files must also exist in the IP pool and must
satisfy the capacity and bandwidth requirements obtained from Level 1.

The three register-file roles can be covered independently. Therefore,
a PE can include all of them or only part of them.

These memories are still kept in the abstract accelerator because
their capacity, bandwidth, and accesses are required by the workload
model. However, their physical implementation is not counted twice.

The \texttt{parent\_idle\_baseline} mode indicates that the idle power
of the included register files is already part of the PE idle power.
Talos only adds their activity-dependent power when the workload is
evaluated.


\subsection{External DRAM}

DRAM is treated as part of the target platform and is not selected by
the Level 2 genome.

A characterized DRAM can be provided in the IP pool. It requires:

\begin{itemize}
    \item bandwidth;
    \item accesses per cycle;
    \item maximum frequency;
    \item idle and active power;
    \item operating-point information.
\end{itemize}

The IP pool can contain at most one DRAM entry. If no DRAM is
provided, Talos creates a fixed platform proxy. In this case, the
characterized IPs in the pool must define one common operating point,
which is also used by the proxy.

The current proxy uses a 512-bit interface, one access per cycle, and
an idle power of $0.02$~W. Its active power is derived from the
Level~1 DRAM energy model.

When a provided DRAM uses \texttt{source: synthetic}, Talos also
derives its active power from the Level~1 model. For a non-synthetic
DRAM characterization, the provided active power is used without
modification.

The automatically generated DRAM model is only a proxy. A real target
platform should replace it with a characterization of the external
memory system.


\subsection{Level 1 Calibration}

Before running the Level 1 exploration, Talos performs the
characterization required by its energy and area models.

For the memory energy model, CACTI is executed using the technology
node specified by the IP pool. This provides the memory access values
used by the Level 1 evaluation.

By default, the relative MAC, register-file, global-buffer, and DRAM
energy costs use the ratios reported for 65~nm:

\[
E_{\mathrm{MAC}} : E_{\mathrm{RF}} : E_{\mathrm{GB}} :
E_{\mathrm{DRAM}} = 1 : 1 : 6 : 200.
\]

These ratios are defined in the Level~1 calibration code and can be
customized when another technology or energy model requires different
relative costs.

The IP pool is also used to obtain the Level 1 area calibration. Talos
looks for compatible physical implementations of the PE, local
storage, and global buffer requirements and uses them to estimate the
area of each abstract architecture.

For this reason, the IP pool must contain enough compatible PE,
register-file, and global-buffer implementations for the architecture
space that is going to be explored.


\subsection{User Constraints}

The user can optionally define physical and workload constraints.

The current implementation supports:

\begin{itemize}
    \item maximum workload latency in cycles;
    \item maximum accelerator area;
    \item maximum average workload power;
    \item minimum implementation frequency.
\end{itemize}

The workload latency constraint can be checked during Level 1.

Area, power, and frequency depend on the selected physical IPs and are
therefore checked during Level 2.

Candidates that do not satisfy the active constraints are removed from
the feasible design space.


\subsection{Exploration Parameters}

Finally, the user can configure how the design space is explored.

The main parameters include:

\begin{itemize}
    \item population size;
    \item number of generations;
    \item number of parallel workers;
    \item random seed;
    \item number of Level 1 architectures passed to Level 2;
    \item Level 2 exploration strategy;
    \item optimization objectives.
\end{itemize}

The standard complete flow uses energy, latency, and area as the
default Level~1 screening objectives. Level~2 uses area, energy, and
workload latency by default. These objective sets belong to the code
that connects and invokes the two levels rather than to the evaluator
implementations themselves. Therefore, a custom flow can call each
level with another combination of the supported objectives.

Level~1 supports latency, energy, area, energy-delay product,
energy-area product, and area-latency product. Level~2 supports area,
energy per inference, average workload power, workload latency,
physical delay, and inverse local-IP throughput.

The default Level~2 strategy is NSGA-II. Exhaustive exploration can be
selected when the remaining physical design space is small enough to
be evaluated completely.

The standard complete flow uses a ZigZag LPF limit of one and
generates one spatial mapping per evaluation. These are default values
of the current calling flow and can be customized when constructing
the ZigZag evaluator or invoking the Level~1 runner from another flow.

When exhaustive exploration is used, a maximum number of combinations
can also be specified to avoid exploring an excessively large physical
design space.

With one objective, exhaustive exploration orders the feasible
solutions directly by that objective. With several objectives, it uses
logarithmic deviations from the ideal and an augmented Tchebycheff
score. The score calculated for one architecture uses a local ideal
and is not comparable with another architecture. The full flow
recalculates a global score using one common ideal over the feasible
solutions being compared.


\subsection{Software Requirements}

Talos is implemented in Python and the current version has been tested
using Python~3.12.

The main external tools and libraries used by the framework are:

\begin{itemize}
    \item ZigZag for workload mapping and architecture evaluation;
    \item pymoo for the NSGA-II explorations;
    \item ONNX for the workload representation;
    \item CACTI for the Level 1 memory characterization.
\end{itemize}

The required Python packages are included in the project dependency
file. CACTI is executed automatically as part of the Level 1
calibration. The source code, the dependency file and a quick guide on
how to install and run the tool (in the repository README) are
available at the Talos repository~\cite{talos_repo}.
